\documentclass[11pt,a4paper]{article}

\usepackage[utf8]{inputenc}
\usepackage[T1]{fontenc}
\usepackage{amsmath,amssymb}
\usepackage{graphicx}
\usepackage{booktabs}
\usepackage{hyperref}
\usepackage{xcolor}
\usepackage{tcolorbox}
\usepackage[margin=25mm]{geometry}
\usepackage{xeCJK}
\setCJKmainfont{Hiragino Mincho ProN}
\setCJKsansfont{Hiragino Sans}

\title{水晶 BAW 実験 初歩からの完全ガイド\\
\large なぜ水晶か，何を測るか，どう違うか}
\author{Wright Brothers Research Program}
\date{2026年4月}

\begin{document}
\maketitle

\begin{abstract}
本ガイドは，水晶バルク音響共振器（Bulk Acoustic Wave resonator,
BAW）を使った Wright Brothers プロジェクトの核心実験を，
前提知識ゼロから理解できるように書かれている．「なぜ水晶なのか」
「なぜ厚さを変えるのか」「なぜ Dirichlet--Neumann 境界にこだわるのか」
「過去の Casimir 実験や水晶実験と何が違うのか」という根本的な
疑問に一つ一つ答えていく．量子力学や場の量子論の教科書的知識は
前提としない（必要な部分は本文で導入する）．実験の各ステップに
ついて「何をするか」だけでなく「なぜそうするか」を重視した．
目標読者は理系大学 1 年生，熱心な高校生，または実験に興味のある
一般の読者．読了後には，¥100 円の時計用水晶振動子で
ダークエネルギーを測りにいく実験を，自分の手で設計できる．
\end{abstract}

\tableofcontents

\newpage
\section{なぜこの実験をやるのか}

まず，この実験が何のためにあるのかを正直に述べる．

我々は\textbf{宇宙の真空エネルギーの謎}を解こうとしている．
現代物理学には「宇宙論定数問題」と呼ばれる巨大な矛盾がある：
理論（量子電磁力学）が予測する真空エネルギー密度と，実際に観測
される値が\textbf{120 桁ずれている}．これは物理学史上最悪の
理論と観測のズレ．

Wright Brothers プロジェクトの仮説：
\begin{itemize}
\item 真空エネルギーは距離に応じて $1/d^4$ で増大するが，
ある臨界距離 $a_c$ でプラトー化する
\item $a_c \approx 88\,\mu$m であり，これはダークエネルギー密度に
対応する自然長スケール
\item この仮説が正しければ，$d = 88\,\mu$m 近辺で \textbf{真空の
振る舞いが変わる}のが直接観測できるはず
\end{itemize}

$88\,\mu$m は髪の毛の太さ程度．そして\textbf{この厚さの水晶板は
100 円で買える}．つまり我々は，100 円の電子部品で宇宙論的謎の
答えを持ってこようとしている．

この実験が本当に成功するかは分からない．しかし，試す価値は明らかに
ある．もし成功すればノーベル賞級，失敗しても物理学の新しい測定
データが得られる．どちらに転んでも損はない．

そして何より：\textbf{この実験は今週発注すれば来月には結果が
出始める}．これが本ガイドのすべてを理解する動機である．

\section{水晶とは何か——物質としての基本}

\subsection{水晶という結晶}

「水晶」はケイ素と酸素からなる結晶の一種で，化学式は
SiO$_2$．地殻の重量の $\sim 12\%$ を占める．ガラスと同じ成分だが，
原子が規則正しく並んでいる点が決定的に違う．

ガラスは原子がランダムに配置された「アモルファス」固体．
水晶は原子が六角形の格子を組む\textbf{単結晶}．この違いが物理
特性を決定する．

\begin{itemize}
\item ガラス: 音が伝わりにくく減衰大（内部摩擦大）
\item 水晶: 音が極めて綺麗に伝わり減衰小（内部摩擦極小）
\end{itemize}

減衰が小さい＝振動が長く続く＝\textbf{機械的 Q 値が高い}．水晶の
Q は室温で $10^6$，液体窒素温度で $10^8$，20\,mK では $10^{10}$ を
超える．これは地上で最も「鳴る」物質の 1 つである．

\subsection{水晶の圧電効果}

水晶の最大の特徴は\textbf{圧電効果}．これは \textbf{1880 年}に
ピエール・キュリーとジャック・キュリー兄弟によって発見された
現象で，以下の 2 つを意味する：

\begin{enumerate}
\item 水晶に力を加えると電気が生じる（直接圧電効果）
\item 水晶に電圧を加えると変形する（逆圧電効果）
\end{enumerate}

これは結晶の\textbf{非中心対称性}から生まれる．水晶の Si と O の
配置には反転対称性がなく，外力で格子がわずかに歪むと正電荷と
負電荷の中心がずれて電気双極子が発生する．逆に，電場をかければ
その双極子が整列しようとして結晶が機械的に変形する．

この性質により，水晶は\textbf{電気信号と機械振動を相互変換}
できる．この変換効率は高くないが（圧電定数 $d_{11} \sim 2$\,pm/V），
それでも十分実用になる．

\subsection{カット方向（orientation）}

水晶は結晶なので，どの方向に切り出すかで性質が変わる．主要な
カット方向と用途：

\begin{center}
\begin{tabular}{lll}
\toprule
カット名 & 角度 & 主用途 \\
\midrule
X カット & 基本軸切り & 歴史的，あまり使われない \\
Y カット & X と垂直 & 温度特性悪い \\
\textbf{AT カット} & $+35.25^\circ$ Y軸 & \textbf{時計・通信の主力} \\
BT カット & $-49^\circ$ Y軸 & 高周波 \\
SC カット & 2 軸回転 & 超高安定（GPS 基準局など）\\
BVA カット & SC 派生 & 研究用超高 Q \\
\bottomrule
\end{tabular}
\end{center}

\textbf{AT カット}が圧倒的に重要．これは Bechmann が 1934 年に
発見したもので，\textbf{室温付近で温度係数がほぼゼロ}になる
特殊な角度．この性質が水晶時計・水晶発振器の大量生産を可能にした．

本実験でも AT カットを使う．$-20^\circ$C から $+60^\circ$C で周波数
変動が数 ppm に抑えられ，実験室環境で扱いやすい．

\subsection{AT カット水晶の音速}

本実験で最も重要な数値：
\begin{center}
\Large
\textbf{AT カット水晶の shear wave 音速 $v_s = 3320$\,m/s}
\end{center}

この 3320\,m/s が，水晶の厚さと共振周波数を結ぶ．厚さ $d$ の
水晶板の基本周波数は $f_1 = v_s/(2d)$ で，例えば $d = 166\,\mu$m
なら $f_1 = 10\,$MHz．

\section{BAW 共振器とは何か}

\subsection{「バルク音響波」の意味}

BAW は ``Bulk Acoustic Wave'' の略．「バルク」は\textbf{表面ではなく
体積内部}を伝わる音響波という意味．対比されるのは SAW
(Surface Acoustic Wave) で，これは表面を伝わる波．

バルクを選ぶ理由：
\begin{itemize}
\item 体積モードは境界の影響を受けにくく Q が高い
\item 高周波化が容易（$f = v/\lambda$，薄くすれば高くなる）
\item 結晶の深い部分の純粋な機械特性を反映する
\end{itemize}

\subsection{厚さモード}

水晶板に電極を挟んで電圧をかけると，逆圧電効果で板が厚さ方向に
伸び縮みする．この伸縮が\textbf{厚さ音響波}として板内部を
上下に往復する．

両面で反射して定常波が立つ条件は，板厚 $d$ が波長の整数倍に
なること：
\begin{equation}
n \cdot \frac{\lambda_n}{2} = d \quad\Rightarrow\quad \lambda_n = \frac{2d}{n}
\end{equation}
周波数に直すと：
\begin{equation}
f_n = \frac{v_s}{\lambda_n} = \frac{n v_s}{2 d}, \quad n = 1, 2, 3, \ldots
\end{equation}

これが両端自由（Neumann-Neumann, NN）の場合．$n=1$ が基本，
$n=2, 3, \ldots$ が倍音（オーバートーン）．

\subsection{なぜオーバートーンが出るか}

物理的には，楽器の弦と同じ．両端固定の弦は基本振動
（全体が一つの山）だけでなく，2 倍振動（2 つの山），3 倍振動
（3 つの山）...も励起できる．BAW でも同じことが起こり，
\textbf{整数倍の周波数にピーク}が並ぶ．

実用の水晶発振器では基本振動だけを使うことが多いが，
本実験では\textbf{オーバートーンの構造そのものが情報}になる．

\subsection{Q 値という指標}

共振器の性能を表す最重要パラメータ：
\begin{equation}
Q = \frac{f}{\Delta f} = 2\pi \cdot \frac{\text{貯蔵エネルギー}}{\text{1 周期あたり散逸}}
\end{equation}
共振ピークの鋭さ．Q が高いほどピークが細く，周波数測定精度が上がる．

\begin{center}
\begin{tabular}{lll}
\toprule
物体 & Q 値 & 備考 \\
\midrule
ゴムボール & 10 & 数回バウンドして止まる \\
ワイングラス & 1,000 & チーン\\
音叉 & 10,000 & 長く鳴る \\
水晶（室温） & $10^6$ & \\
水晶（液体窒素）& $10^8$ & 77\,K \\
水晶（希釈冷凍機）& $10^{10}$ & 20\,mK, BVA カット \\
超伝導マイクロ波共振器 & $10^{11}$ & \\
光学共振器（高級品）& $10^{11}$ & 超精密 FP \\
\bottomrule
\end{tabular}
\end{center}

水晶の $10^{10}$ というのは機械振動系として世界最高クラス．
これが本実験の精度を保証する．

\section{水晶と精密時間の歴史}

水晶 BAW 実験を理解するには，水晶がどうやって人類最高の時間
基準になったか知っておくと良い．

\subsection{1920 年代：Marrison の発明}

1927 年，Warren Marrison（ベル研）が世界初の水晶時計を作成．
室温での精度は 1 日に約 0.01 秒，当時の最高性能を凌駕した．
振り子時計と比較した優位性は：
\begin{itemize}
\item 機械振動ではなく電気的な発振器
\item 外部振動の影響を受けにくい
\item 小型化・量産化が容易
\end{itemize}

\subsection{1930 年代：AT カットの発見}

1934 年，Bechmann が AT カットを発見．温度係数が $35.25^\circ$ で
ほぼゼロになる．これにより温度変動環境でも安定した発振器が
可能になり，\textbf{通信・レーダー・時計の基礎技術}となった．

\subsection{1960 年代：原子時計との比較}

セシウム原子時計（1955 年）が登場したことで「1 秒」の定義が
天文学から原子物理に移行．しかし水晶発振器は\textbf{原子時計の
短期安定度の足場}として依然重要．GPS 衛星の内部時計も
水晶とセシウムのハイブリッド．

\subsection{1990 年代：BVA カットと超高 Q}

Besson が 1976 年に開発した BVA (Boîtier à Vieillissement
Amélioré) カットは，エッジ損失を最小化する特殊な保持方式．
SC カットの一種で，室温で $Q > 10^6$，低温で $10^8$ 以上に達する．

このカットが後に Galliou, Goryachev, Tobar らのグループによって
\textbf{20\,mK まで冷やされ，$Q > 10^{10}$ を達成}．これは
機械振動系で人類が到達した最高の Q 値である．

\subsection{本実験との関係}

本実験は，この 100 年にわたる水晶技術の蓄積に乗っかる．特に：
\begin{itemize}
\item AT カット $\to$ 温度安定性
\item 量産品 $\to$ 安価 (¥100 円)
\item 厚さの範囲 $\to$ 自然に $a_c = 88\,\mu$m を含む
\item 高 Q $\to$ 微小な周波数シフトの検出が可能
\end{itemize}

水晶以外の材料ではこの組み合わせが揃わない．例えば LiNbO$_3$
(ニオブ酸リチウム)は Q が低く，AlN (窒化アルミ)は加工が難しく，
ZnO (酸化亜鉛)は安定性が劣る．水晶一択．

\section{境界条件という概念}

ここからが物理的な核心．境界条件とは何か，なぜ重要かを理解する．

\subsection{弦と管の基礎}

まず音楽の例から．ギターの弦は両端が固定（クランプ）されている．
このとき定常波の条件は「両端で振幅ゼロ」．これを\textbf{ディリクレ
境界条件}と呼ぶ．数式で書くと $u(0) = u(L) = 0$．

一方，パイプオルガンやフルートの管は一方または両方が開いている．
開放端では空気が自由に出入りでき，\textbf{圧力}がゼロ（大気圧に
等しい）．これを\textbf{ノイマン境界条件}と呼ぶ．数式で書くと
$\partial u/\partial x|_{\text{開端}} = 0$．

\subsection{4 通りの境界条件}

境界条件の組み合わせは 4 通り：

\begin{center}
\begin{tabular}{llll}
\toprule
名前 & 略記 & 両端 & 例 \\
\midrule
Dirichlet--Dirichlet & DD & 両端固定 & ギター弦，両端閉管 \\
Neumann--Neumann & NN & 両端自由 & 両端開放管，宙吊り板 \\
Dirichlet--Neumann & DN & 片端固定片端自由 & 片端閉管，片面クランプ板 \\
Neumann--Dirichlet & ND & DN と等価（向き逆）& 同上 \\
\bottomrule
\end{tabular}
\end{center}

DD と NN は同じスペクトル（どちらも全倍音 $f_n = n f_0$）．
DN は\textbf{奇数倍音のみ} $f_n = (2n-1)f_0$．これが決定的に重要．

\subsection{なぜ DN は奇数倍音だけか}

物理的に直感的に理解するには，ハーフサイズで考える．長さ $2L$ の
両端自由管は基本波が $\lambda = 2L$ で $n=1,2,3,\ldots$ の全倍音．
この管の真ん中で切ると，左半分は「左端自由，右端ずっと固定
（節）」となる．これは長さ $L$ の DN 管．

元の $2L$ 管で真ん中に節を持てる倍音は $n=2, 4, 6,\ldots$（偶数）
のみ．奇数倍音は真ん中が腹になる．したがって「真ん中で切る」操作は
偶数倍音を切り捨てる操作と同じ．結果として $L$ の DN 管は奇数倍音
のみ．

逆に言えば，\textbf{DN 管は同じ長さの DD/NN 管から偶数倍音を
取り除いたもの}．

\subsection{なぜ境界条件が物理的に重要か}

真空は量子力学的には「全モードの基底状態」である．各モードは
零点エネルギー $\hbar\omega/2$ を持つ．総和は：
\begin{equation}
E_{\rm vac} = \sum_n \frac{\hbar\omega_n}{2}
\end{equation}

NN と DN でモードの集合が違うので，真空エネルギーが異なる．
その差が物理的に観測可能なのが Casimir 効果．

より抽象的には：\textbf{「真空」とは境界条件に依存する}．無限に
広い真空と，境界を持つ真空は別物．本実験はこの違いを定量化する．

\section{Dirichlet--Neumann 境界の具体化}

\subsection{機械的な DN の作り方}

電磁場の DN 境界は実装困難（完全磁気導体 PMC が必要だが広帯域で
存在しない）．しかし\textbf{機械系の DN は簡単}：

\begin{itemize}
\item 自由表面（空気 or 真空に接する）: stress $= 0$ なので Neumann
\item クランプ表面（重く動かない基板に接着）: displacement $= 0$ なので Dirichlet
\end{itemize}

したがって水晶板の片面をタングステン板に接着するだけで DN が完成．

\subsection{なぜタングステンか}

音響インピーダンス $Z = \rho v$（密度×音速）の大きさで決まる．

\begin{center}
\begin{tabular}{lll}
\toprule
材料 & $\rho$ (g/cm³) & $Z$ (MRayl) \\
\midrule
空気 & 0.0012 & $0.0004$ \\
水 & 1.0 & 1.5 \\
アルミニウム & 2.7 & 17 \\
\textbf{水晶} & \textbf{2.65} & \textbf{15} \\
銅 & 8.9 & 42 \\
サファイア & 4.0 & 44 \\
\textbf{タングステン} & \textbf{19.3} & \textbf{100} \\
\bottomrule
\end{tabular}
\end{center}

水晶とタングステンのインピーダンス比 $100/15 = 6.7$．これが
音響反射の強さを決める．反射率は：
\begin{equation}
R = \left|\frac{Z_2 - Z_1}{Z_2 + Z_1}\right|^2 = \left|\frac{100-15}{100+15}\right|^2 = 0.547
\end{equation}

54\% 反射＝46\% 透過．これではまだ Dirichlet 近似として弱い．

\subsection{反射率を上げる工夫}

方法 1: 多層化．タングステン 5\,mm だけだとその裏も境界があるので，
タングステン－水－タングステン－ などの多重反射で実効的に
$>90\%$ 反射を達成．

方法 2: 音響 Bragg 反射鏡（SMR が使っている方法）．$\lambda/4$
厚の高インピーダンス/低インピーダンス層を交互に積層．
詳細は \texttt{guide\_smr\_fbar\_ja.tex}．

方法 3: タングステンの代わりに白金 ($Z = 75$) や金 ($Z = 63$)．
安定した接着が可能．

本実験では\textbf{タングステン 10\,mm 厚単体}で十分．反射率は
$55\%$ だが，系統誤差として処理できる．完璧な Dirichlet でなくとも
「$1/d^4$ からの逸脱」という定性的な signature は見える．

\subsection{接着剤の選択}

接着剤自体が音響的に中間層を作る．理想はゼロ厚さ．

\begin{itemize}
\item \textbf{Stycast 1266}（低温用エポキシ）: 1--5\,$\mu$m 厚，
音響特性既知
\item \textbf{Indium 溶着}: 金属接合，理想的だが加工困難
\item \textbf{シアノアクリレート}: 薄く塗布可能だが低温で剥離
\end{itemize}

Wright Brothers Phase 2 では Stycast 1266 を使う．
厚さを制御するためにガラスビーズ（$5\,\mu$m 径）を混ぜて
スペーサとする．

\section{実験の全体像}

ここまでで基礎知識が揃った．実験全体を俯瞰する．

\subsection{目的}

厚さ $d$ の異なる水晶 BAW を複数用意し，それぞれを NN (両面自由)
と DN (片面クランプ) の 2 形態で測定する．両者の共振周波数の
差$\Delta f(d)$ を $d$ の関数としてプロットし，$1/d^n$ の
スケーリング指数 $n$ を決定する．

\begin{itemize}
\item $n \approx 3$ なら標準 QED 支持
\item $n < 3$ なら新物理の兆候
\item $n$ が変化する点があればプラトー発見
\end{itemize}

\subsection{なぜ厚さを変えるのか}

一つの厚さでは「$1/d^n$ の $n$」が決まらない．複数の $d$ で
測定して対数プロットを取り，傾きを抽出する．

さらに，我々の予測する $a_c \approx 88\,\mu$m は，ある特定の
厚さで何かが起こるという主張．\textbf{その厚さ付近でスペクトル
振る舞いが変わる}のを見るには，その厚さを含む範囲を網羅する
必要がある．

\subsection{使う厚さ範囲}

Phase 2 の設計：

\begin{center}
\begin{tabular}{lll}
\toprule
厚さ $d$ & 基本周波数 $f_1$ & 位置付け \\
\midrule
830\,$\mu$m & 2\,MHz & 下限（プラトー内？）\\
330\,$\mu$m & 5\,MHz & \\
166\,$\mu$m & 10\,MHz & \\
\textbf{83\,$\mu$m} & \textbf{20\,MHz} & \textbf{$a_c$ 近傍}\\
33\,$\mu$m & 50\,MHz & \\
17\,$\mu$m & 100\,MHz & 上限 \\
\bottomrule
\end{tabular}
\end{center}

50 倍のレンジで対数的に 6 点．これで $\log d$ vs $\log\Delta f$
プロットに意味のあるフィットが可能．

\subsection{入手方法}

全て\textbf{商用水晶振動子}から取れる．秋月電子，Digi-Key，
RS オンラインなどで 1 個 ¥100--¥1000．

\begin{itemize}
\item 時計用 32.768 kHz はチューニングフォーク型で不適
\item 発振器用の円板型を選ぶ
\item 「HC-49」「HC-51」「FT-23」パッケージのものが standard
\end{itemize}

パッケージを割って中の水晶円板を取り出す作業が必要（缶オープナー
とラジオペンチで可能）．

\section{ステップごとの意味}

いよいよ実験の具体的な各ステップを，「何をして」「なぜそうするか」
の両方から説明する．

\subsection{Step 1: 水晶の入手と準備}

\textbf{何をするか}: 各周波数の水晶振動子を秋月等で購入．
パッケージを開けて中の水晶を取り出す．

\textbf{なぜか}: 商用品はメタルパッケージ内に入っているが，
\textbf{パッケージ自体が機械共振系}なので，直接水晶を取り出して
自前のマウントに移す必要がある．メタル缶のまま測定すると缶の
モードと混ざって解釈不能．

\textbf{注意点}:
\begin{itemize}
\item 水晶は非常に割れやすい．工具の先端を柔らかい布で巻く
\item 電極（金蒸着）を指で触ると汚染する．ピンセット必須
\item 水晶に残った接着剤をアセトンで洗浄
\end{itemize}

\subsection{Step 2: NN 試料の作成}

\textbf{何をするか}: 取り出した水晶を，両面が空気に触れる形で
マウントする．通常は 3 点支持（エッジを 3 つの細いサポートで
保持）．

\textbf{なぜか}: 両面自由＝NN 境界．これが基準線．3 点支持は
マウンティング損失を最小化する古典的方法（Besson 1976）．

\textbf{必要なもの}:
\begin{itemize}
\item サポートピン（細い銅線 $\phi 0.3$\,mm）
\item 真空対応セラミック基板
\item 電極接続用の導電性エポキシ（最小量）
\end{itemize}

\subsection{Step 3: DN 試料の作成}

\textbf{何をするか}: 別の水晶（同じ厚さ）の片面をタングステン板に
接着する．接着剤は Stycast 1266 を使い，ガラスビーズで
厚さを制御．

\textbf{なぜか}: 片面をクランプ＝Dirichlet，もう片面は自由
＝Neumann．これで DN 境界が完成する．タングステンが十分重く
動かないのでクランプ近似が成立．

\textbf{手順}:
\begin{enumerate}
\item タングステン板 ($5\,\text{mm} \times 20\text{mm} \times 20\text{mm}$) を
研磨
\item 脱脂（アセトン，イソプロパノール）
\item Stycast 1266 に $5\,\mu$m ガラスビーズを $1\%$ 混合
\item 水晶片面にミクロスポイトで少量塗布
\item タングステン板に接着，3 時間室温硬化
\item 反対面（自由側）は空気
\end{enumerate}

\textbf{注意}: 接着剤の厚さと均一性が境界条件の品質を決める．
ビーズがスペーサとなり $5\,\mu$m に揃う．

\subsection{Step 4: 電極の取り付け}

\textbf{何をするか}: 水晶の両面（または DN 側の自由面と
タングステン板）に電気的接続を取る．

\textbf{なぜか}: 圧電効果で振動を励起・検出するため．水晶の両面に
薄い金電極がすでに蒸着されているのが通常．これに細いワイヤを
接続．

\textbf{方法}:
\begin{itemize}
\item 導電性エポキシ（Epo-Tek H20E）をごく少量
\item ワイヤは $\phi 25\,\mu$m の金線または $\phi 50\,\mu$m 銅線
\item 圧電振動を妨げないよう接続点は小さく
\end{itemize}

\subsection{Step 5: 真空封入}

\textbf{何をするか}: 試料を真空容器に入れ，$< 10$\,Pa まで排気．

\textbf{なぜか}: 空気の粘性減衰は Q を大きく下げる．室温で
$10^{-3}$ Torr 以下なら Q は 10 倍以上改善する．\textbf{高 Q が
微小な周波数シフトの測定精度を決める}．

\textbf{簡易版}:
\begin{itemize}
\item デシケータ改造（¥10,000）
\item 真空ポンプ（中古ロータリ $\sim$¥30,000）
\item バキュームゲージ $\sim$¥5,000
\end{itemize}

\subsection{Step 6: 温度安定化}

\textbf{何をするか}: 試料の温度を $\pm 0.1\,$K 以内に制御．

\textbf{なぜか}: 水晶の寸法は熱膨張し，共振周波数が温度でドリフト
する．AT カットは温度係数が小さいが，精密測定では無視できない．
$25^\circ\text{C} \pm 0.1\,\text{K}$ 程度で安定化すれば十分．

\textbf{方法}:
\begin{itemize}
\item ペルチェ素子 + アルミ熱拡散ブロック
\item PID 温度コントローラ（Arduino ベース）
\item 熱電対 or サーミスタで温度モニタ
\end{itemize}

\subsection{Step 7: 共振スペクトル測定}

\textbf{何をするか}: nanoVNA または専用計測器で，試料の電気的
インピーダンス $Z(f)$ を周波数掃引で測定．

\textbf{なぜか}: 水晶の機械共振は，電気的には\textbf{インピーダンスの
極小}として現れる（圧電効果の逆変換）．この極小の位置が共振
周波数 $f_1$，幅が $Q$ を教える．

\textbf{実際の操作}:
\begin{itemize}
\item nanoVNA を SOL 校正
\item 試料を接続
\item 基本周波数近傍を狭帯域スイープ（$\pm 10\%$）
\item ローレンツ関数でフィット
\item $f_1$ と $Q$ を抽出
\end{itemize}

\subsection{Step 8: オーバートーン測定}

\textbf{何をするか}: 基本だけでなく $2f_1, 3f_1, 4f_1, \ldots$ まで
測定．

\textbf{なぜか}: DN では偶数倍音が欠落しているはず．その欠落の
確認が\textbf{境界条件の物理的実現の証拠}．NN 試料と DN 試料で
オーバートーン構造が本当に違うかを見る．

\textbf{重要}: 単に「ピークがある／ない」だけでなく，$Q$ も
測定．DN の偶数倍音位置では，もし完全に欠落していれば Q は
測定不能（ピーク自体がない）．不完全な DN なら減衰した
ピークが見える．その度合いが Bragg 反射鏡の品質を示す．

\subsection{Step 9: 非線形駆動応答}

\textbf{何をするか}: 駆動パワーを $-20$\,dBm から $+5$\,dBm まで
変えて，共振周波数の変化を測定．

\textbf{なぜか}: 強く駆動すると非線形効果（Duffing 型）で周波数が
シフトする．この シフト量は高次モードとの結合 $g_n$ の 2 乗に
比例．NN と DN で $g_n$ の集合が違うので，シフトの係数も違う．
\textbf{これが Lamb shift の代理観測}．

\subsection{Step 10: 温度スイープと差分解析}

\textbf{何をするか}: 温度を $20^\circ$C から $50^\circ$C まで $1^\circ$C
刻みで変え，各温度で共振周波数を測定．$df/dT$ を取る．

\textbf{なぜか}: 加工誤差や個体差による絶対周波数のズレを
キャンセルするため．温度依存性は材料固有のはずなので，
NN と DN で同じ温度係数になるべき．\textbf{異なる温度係数が
モード構造由来の補正}を示す．

これは Phase 2 の\textbf{最も重要な測定}．ここで 10\,mHz オーダーの
Lamb shift 差が検出できる．

\subsection{Step 11: 厚さシリーズ解析}

\textbf{何をするか}: Step 1--10 を全 6 厚さで繰り返す．各厚さで
$\Delta f = f_{NN} - f_{DN}$（または非線形シフト差）を得る．
$\log d$ vs $\log |\Delta f|$ プロットを作り傾きをフィット．

\textbf{なぜか}: これが\textbf{本実験の本丸}．傾きが $-3$ なら
標準 QED，それより浅ければ新物理，変化点があればプラトー．
$a_c = 88\,\mu$m に対応する厚さ近傍での振る舞いが決定的．

\subsection{Step 12: 論文化}

\textbf{何をするか}: Phase 2 完了までに Paper 2 or 3 のドラフト
を進め，arXiv にプレプリント投稿．

\textbf{なぜか}: 科学コミュニティに結果を公開することで，
\begin{itemize}
\item フィードバックが得られる
\item 優先権が確保される
\item 次段階（Phase 3, 4）の共同研究交渉材料になる
\end{itemize}

\section{過去の実験との違い}

ここで，なぜ今までこの実験が行われていないのかを整理する．

\subsection{従来の Casimir 実験との違い}

\begin{center}
\begin{tabular}{lll}
\toprule
比較項目 & 従来 & 本実験 \\
\midrule
境界 & DD（金属ー金属）& DN（水晶自由ークランプ）\\
距離 & nm--$\mu$m & $\mu$m--mm \\
測定量 & 力（AFM, トーション）& 周波数 \\
系 & 電磁真空 & 機械的フォノン真空 \\
Q 値 & n/a & $10^6$ 以上 \\
プラトー検証 & 単一距離 & 厚さシリーズ \\
\bottomrule
\end{tabular}
\end{center}

\textbf{本質的な新規性}：
\begin{enumerate}
\item DN 境界を真面目に実装
\item 距離ではなく厚さを体系的に変える
\item 力ではなく周波数で測る（桁違いに精度高い）
\item 機械系を使うので Kramers--Kronig 制約がない
\end{enumerate}

\subsection{Lamoreaux 1997 (最初の精密 Casimir)}

\begin{itemize}
\item トーション振り子と球－平板
\item DD 境界のみ
\item 距離 $0.6$--$6\,\mu$m
\item $1/a^4$ 則を 5\% 精度で確認
\item 本実験は\textbf{別の境界条件・別の距離領域}
\end{itemize}

\subsection{Galliou, Goryachev, Tobar (超高 Q 水晶)}

\begin{itemize}
\item BVA カット水晶，20\,mK
\item $Q > 10^{10}$ 達成
\item 重力波・ダークマター検出器として使用
\item \textbf{境界は NN のみ}．DN 構造との比較は未実施
\item 本実験は\textbf{彼らの手法を DN で使う}ことの提案
\end{itemize}

\subsection{Chu, Schoelkopf (HBAR 量子音響)}

\begin{itemize}
\item 薄膜 AlN/サファイア HBAR
\item 単一音子状態生成・操作
\item \textbf{両面自由（NN）構造}
\item 単一モードのみ扱い，スペクトル comb 全体は見ない
\item 本実験は\textbf{異なる境界と異なる観点}
\end{itemize}

\subsection{Adelberger (Eöt-Wash 重力)}

\begin{itemize}
\item トーション天秤で重力逆二乗則テスト
\item $\sim 50\,\mu$m まで null 結果
\item \textbf{通常物質（タングステン）のみ，DN 境界なし}
\item 本実験は\textbf{境界を持つ系で初めて $50\,\mu$m 領域に踏み込む}
\end{itemize}

\subsection{Boyer (1974)}

\begin{itemize}
\item DN 境界で Casimir 斥力を\textbf{理論的に}予測
\item 球体での計算：$E_{\rm Boyer} = +0.0916 \hbar c/r$
\item \textbf{平板 DN の実験は 50 年間誰もやっていない}
\item 本実験は\textbf{Boyer の予言の初めての直接検証}
\end{itemize}

\subsection{なぜ誰もやっていないのか}

正直な分析：

\begin{enumerate}
\item \textbf{DN 境界が難しいという誤解}：電磁版は確かに難しいが
機械版（水晶＋重基板）は簡単
\item \textbf{分野の分断}：Casimir は電磁，BAW は RF 工学，
$K$ 理論は純粋数学．三者をつなぐ発想がなかった
\item \textbf{厚さシリーズという発想}：Casimir は距離スキャンが
主流で，厚さは固定されがち
\item \textbf{プラトー予想自体が新しい}：Wright Brothers 固有の
理論的動機づけがなければ「なぜ厚さを変える必要があるか」が
分からない
\end{enumerate}

つまり：\textbf{やれる技術は 50 年前からあった．発想がなかっただけ}．

\section{何が起きうるか：予想されるシナリオ}

実験が進むにつれて起きうる事態を分類しておく．

\subsection{シナリオ A: 標準 QED 通り（確率 $\sim 40\%$）}

$\log d$--$\log |\Delta f|$ プロットで傾きが全域で $-3$．
これは現代物理の予測通り．つまらない結果だが価値はある：

\begin{itemize}
\item DN レジームの初の系統測定
\item プラトーが存在するなら $d_{\min} < 17\,\mu$m 以下
\item 論文 2, 3 として publish 可能
\item 次は薄膜 SMR（mmWave）でさらに薄い側を探索
\end{itemize}

\subsection{シナリオ B: 傾きが浅い（確率 $\sim 30\%$）}

傾きが $-2$ 程度．プラトー手前か，別の新物理．

\begin{itemize}
\item 最も興味深い中間状態
\item 追加測定（低温，精密化）が必要
\item Phase 3 分岐 A へ
\end{itemize}

\subsection{シナリオ C: プラトー発見（確率 $\sim 20\%$）}

$d \sim 88\,\mu$m 近傍で傾きが急変．$d > a_c$ でほぼ flat．

\begin{itemize}
\item 革命的結果
\item ダークエネルギー問題の実験的解決候補
\item 即時論文化（PRL → Nature Physics）
\item Phase 3 分岐 A で精密化
\end{itemize}

\subsection{シナリオ D: 偶数倍音欠落が見えない（確率 $\sim 10\%$）}

DN 試料でも NN と同じスペクトルが出る．

\begin{itemize}
\item DN 境界の実現が不完全
\item タングステン反射不足
\item 接着剤厚さの問題
\item Phase 3 分岐 C（系統改善）
\end{itemize}

確率はあくまで主観的評価．理論的には C の方が物理的に豊かな帰結を
もたらすので，優先度高く監視する．

\section{落とし穴と対策}

\subsection{接着剤の問題}

\textbf{問題}: 接着剤層が厚すぎるとクランプ近似が崩れる．薄すぎると
強度不足．

\textbf{対策}: ガラスビーズスペーサで $5\,\mu$m に制御．接着後に
光学顕微鏡で断面観察．

\subsection{電極の質量負荷}

\textbf{問題}: 水晶の両面の金電極（通常 $100$\,nm 厚）は質量負荷と
して働き，共振周波数を下げる．

\textbf{対策}: 両試料（NN, DN）で電極厚を揃える．もし可能なら
同じロットから取る．質量効果は NN/DN で共通なので差分では消える．

\subsection{スプリアスモード}

\textbf{問題}: 厚さ方向以外の横モード（Lamb 波，たわみモード）が
混じる．

\textbf{対策}:
\begin{itemize}
\item 円板状の試料でエネルギーを中心に閉じ込める
\item サポートはエッジのみ，中央は触れない
\item モード同定：加振点を変えて応答パターンから判別
\item $Q$ でフィルタ：スプリアスは $Q$ が低い
\end{itemize}

\subsection{温度ドリフト}

\textbf{問題}: 周波数測定中に温度がドリフトすると絶対値が変わる．

\textbf{対策}:
\begin{itemize}
\item 測定前に 30 分以上温度安定化
\item 連続測定は短時間（$<$5 分）
\item 長時間測定は複数短時間の平均
\item 基準水晶発振器と同時測定で相対ドリフト
\end{itemize}

\subsection{電磁干渉}

\textbf{問題}: 50\,Hz 電源ハムや WiFi (2.4 GHz) がスペクトルに
入る．

\textbf{対策}:
\begin{itemize}
\item 金属シールドボックス
\item バッテリ駆動
\item 測定中は WiFi/Bluetooth 切
\item ツイストペアケーブル
\end{itemize}

\subsection{サンプル破損}

\textbf{問題}: $17\,\mu$m の薄い水晶は非常に壊れやすい．

\textbf{対策}: スペア複数個．取り扱いは真空ピンセット．絶対に指で
触らない．

\section{最小予算での現実的開始プラン}

ここで「本当に始めるとしたら」の具体的予算を提示する．

\subsection{第 0 週: 発注}

\begin{center}
\begin{tabular}{lll}
\toprule
項目 & 単価×数 & 小計 \\
\midrule
時計用水晶 2 MHz ×2 & ¥500 & ¥1,000 \\
時計用水晶 5 MHz ×2 & ¥100 & ¥200 \\
時計用水晶 10 MHz ×2 & ¥100 & ¥200 \\
時計用水晶 20 MHz ×2 & ¥200 & ¥400 \\
時計用水晶 50 MHz ×2 & ¥500 & ¥1,000 \\
時計用水晶 100 MHz ×2 & ¥1,000 & ¥2,000 \\
タングステン板 $5\text{mm}\times 20\times 20$ ×6 & ¥1,000 & ¥6,000 \\
Stycast 1266 エポキシ & & ¥8,000 \\
ガラスビーズ $5\,\mu$m & & ¥3,000 \\
導電性エポキシ Epo-Tek H20E & & ¥5,000 \\
金線 $\phi 25\,\mu$m 1\,m & & ¥4,000 \\
ピンセット（精密）& & ¥2,000 \\
セラミック基板 & & ¥3,000 \\
nanoVNA V2 Plus4 & & ¥25,000 \\
周波数カウンタ中古 & & ¥60,000 \\
温度コントローラ DIY & & ¥15,000 \\
デシケータ & & ¥8,000 \\
真空ポンプ中古 & & ¥30,000 \\
その他 工具・配線 & & ¥20,000 \\
\midrule
\textbf{合計} & & \textbf{¥193,800} \\
\bottomrule
\end{tabular}
\end{center}

約 ¥20 万円．1 ヶ月以内に全て揃う．

\subsection{第 1 ヶ月: 準備}

\begin{itemize}
\item Week 1: 水晶取り出し，準備
\item Week 2: NN/DN 試料作成（各厚さ 2 個）
\item Week 3: 電極，マウント
\item Week 4: 真空封入，初期校正
\end{itemize}

\subsection{第 2--3 ヶ月: 測定}

\begin{itemize}
\item 各試料の共振スペクトル
\item オーバートーン解析
\item 非線形駆動応答
\item 温度スイープ差分
\end{itemize}

\subsection{第 4 ヶ月: 解析と論文化}

\begin{itemize}
\item 厚さシリーズ プロット
\item 傾きフィット
\item 不確かさ解析
\item arXiv ドラフト
\end{itemize}

\section{よくある質問 (FAQ)}

\paragraph{Q1: 100 円の水晶で本当にノーベル賞級の実験ができるの？}

原理的には YES．精度と実験手法の問題で，部品の値段ではない．
NASA の一部機器も安価な市販品を使っている．重要なのは
「何を測るか」と「どう解析するか」．

\paragraph{Q2: 個人がこれをやって意味あるの？}

意味はある．Wright Brothers プロジェクトはまさに個人研究者による
貢献を目指している．現代物理でも，Adelberger は大学の
トーション天秤で新発見をしている．\textbf{設備より発想の方が
希少資源}．

\paragraph{Q3: なぜプロの研究者がやっていないの？}

分野の分断（Casimir・BAW・数論）と発想の欠如．誰かがやるべき
だが誰もやっていない領域の典型．Wright Brothers が最初に
飛び込む理由はここにある．

\paragraph{Q4: 失敗したら何が残る？}

\begin{itemize}
\item DN 境界での系統測定データ（新規）
\item 機械 Lamb shift の精密測定技術
\item 水晶 BAW の非標準応用論文
\item Phase 3--5 へ進む土台
\end{itemize}

失敗は「プラトー発見できなかった」というだけ．実験データは
必ず残り，論文化可能．

\paragraph{Q5: 怪我する危険は？}

機械工具と電気工作の一般的注意を守れば安全．エポキシは換気．
真空ポンプの使用は取説通り．放射性物質も高圧高電圧もない．

\paragraph{Q6: 必要な知識レベルは？}

高校物理（波の基礎，電気回路）と，Python 等でのデータ解析の基礎．
本ガイド＋標準的な電子工作入門書で十分．専門的な量子力学は
不要（ただし結果を論文化する段階では役立つ）．

\paragraph{Q7: 時間はどれくらいかかる？}

Phase 2 全体で 3--4 ヶ月．週末に集中できれば可能．平日の仕事が
ある場合は 6 ヶ月ペース．途中途中で何か見えてくるので，
モチベーションは維持しやすい．

\paragraph{Q8: 一人で大丈夫？}

可能．ただし以下の段階では協力者がいると良い：
\begin{itemize}
\item 機械工作（タングステン切削）→ 町工場に依頼
\item 電子回路（ロックイン系）→ 大学の電子工作サークル
\item データ解析・論文書き → 物理学科の大学院生
\end{itemize}

\section{理論的背景の最低限}

実験を理解する上で知っておきたい理論的概念の最短まとめ．

\subsection{ゼロ点エネルギー}

量子力学では，基底状態（最低エネルギー状態）でも有限のエネルギーを
持つ．調和振動子の基底エネルギーは $\hbar\omega/2$．完全にゼロに
なることはない．

これが「揺らぎ」と呼ばれる現象で，真空中でもエネルギーが存在する
原因．

\subsection{Casimir 効果}

2 枚の導体板を近づけると引力が働く．これは板の間の電磁場モードが
離散化（量子化）され，板の外の連続モードとのエネルギー差が板を
引き寄せる．1948 年に Casimir が予言，1997 年に Lamoreaux が
精密検証．

基本公式：
\begin{equation}
F_{\rm Cas} = -\frac{\pi^2 \hbar c}{240 a^4} \cdot A
\end{equation}

\subsection{Boyer の反転}

1974 年，Boyer は「完全電気導体と完全磁気導体の組み合わせ」
（DN 境界）では Casimir 力の符号が反転することを予言．斥力．
現実には PMC が存在しないため電磁での実現困難．

機械系ではクランプー自由で同等の DN が作れる．\textbf{本実験の
中心予言はこれの定量化}．

\subsection{Euler 積}

Riemann ゼータ関数：
\begin{equation}
\zeta(s) = \sum_{n=1}^\infty \frac{1}{n^s} = \prod_p \frac{1}{1 - p^{-s}}
\end{equation}
この無限積（Euler 積）から素数 $p=2$ の因子を除いたもの：
\begin{equation}
\zeta_{\neg 2}(s) = (1 - 2^{-s})\cdot\zeta(s)
\end{equation}
$s = -3$ で $\zeta(-3) = 1/120$, $\zeta_{\neg 2}(-3) = -7/120$．
符号反転と 7 倍の強化．

\subsection{Wright Brothers 仮説}

「$\zeta$ から Euler 積の $p=2$ 因子を切除する操作（$\zeta \to
\zeta_{\neg 2}$）は，境界条件 DD $\to$ DN と同型である」．

さらに「$K_1(\mathbb{Z}[1/2])$ の位相的 winding number が真空
エネルギーをプラトー化する」．

詳細な定式化は \texttt{topological\_casimir.tex},
\texttt{boyer\_reinterpretation.tex} 参照．

\subsection{ダークエネルギーとの関係}

観測される宇宙のダークエネルギー密度 $\rho_\Lambda \approx
5.24\times 10^{-10}$ J/m³ に対応する長さは $\hbar c /\rho_\Lambda^{1/4}
\approx 88\,\mu$m．この長さが\textbf{水晶の標準的な厚さの範囲の
真ん中}にあるのが Wright Brothers 仮説の中心的な動機．

\section{次のステップ}

Phase 2（本実験）が完了したら：

\begin{enumerate}
\item 結果を arXiv プレプリントに（\textbf{Paper 3}）
\item Phase 3（低温化 or SMR/FBAR 精密化）に進む
\item Phase 4 共同研究交渉開始（Galliou, Chu, 村田など）
\item 理論側の強化（$K_1$ regulator 計算を理論家に依頼）
\end{enumerate}

\section{読者への最後のメッセージ}

本ガイドの目的は，「物理の最先端は高額な装置の独占物ではない」
という事実を実感してもらうことだった．水晶は 100 年前から
量産されていて，物理は 50 年前から知られている．必要なのは
「それを試したことがない」という空白に気づき，実際に試すこと．

Wright Brothers の兄弟（Wilbur と Orville）は自転車屋だった．
彼らは当時の航空学の専門家ではなく，しかし初飛行を成功させた．
その理由は，既存の理論を鵜呑みにせず，自分の手で実験を積み重ね，
一つ一つの仮説を検証していったからだ．

本プロジェクトの名前はここから取られている．

\textbf{水晶が届いたら，まず手に取ってみてほしい}．指で弾くと
微かに鳴る音が聞こえるかもしれない．その振動の中に，もしかしたら
宇宙の謎の答えが隠れている．

始めよう．

\vspace{2em}
\hrule
\vspace{0.5em}
\small
\noindent
関連ドキュメント：
\begin{itemize}
\item \texttt{master\_experiment\_plan\_ja.tex}: 全体ロードマップ
\item \texttt{guide\_baw\_quartz\_ja.tex}: 技術的詳細
\item \texttt{guide\_smr\_fbar\_ja.tex}: 商用品ルート
\item \texttt{dark\_energy\_plateau\_ja.tex}: 理論的動機
\item \texttt{boyer\_reinterpretation\_ja.tex}: 数論的再解釈
\end{itemize}

\end{document}
